\section{Light-curve processing} \label{app:processing}

\subsection{Cleaning and detrending}

Cleaning and detrending are unchanged from draft v2, except for the ephemeris used for masking. The light curve is
cleaned with the quality mask of its product. It is then detrended by \code{lcio.detrend} from the frozen helper module
\code{deps/lcio.py}, which uses the biweight filter of wotan \citep{Hippke2019} with a window $W = 0.75$\,d. When the
trend is fitted, both signals are masked over $\pm 1.5\,\Tfourteen$ around their predicted mid-times. Upward outliers
above 4 robust standard deviations ($4 \times 1.4826 \times$ the median absolute deviation) are removed.

\subsection{Segments and edge zones} \label{app:edges}

\code{lcio.detrend} gives wotan the cleaned light curve without the masked windows of both signals ($\pm
1.5\,\Tfourteen$), with a break tolerance of 0.3\,d. Wotan therefore fits a separate segment wherever this series has a
gap longer than 0.3\,d. The edge zones are taken from exactly these segments, for the real detrending and separately
for every fake detrending of the noise model (Appendix~\ref{app:noise}). The first and last $W/2$ of every segment are
edge zones: 0.375\,d (9\,h) for the primary run with $W = 0.75$\,d, and 0.25\,d and 0.5\,d for variants (a) and (b),
whose windows are 0.5\,d and 1.0\,d.

The reason is that the filter is centred and $W$ wide. Within $W/2$ of a segment end it has less than half a window of
data on one side, so its trend there is less reliable. The wotan documentation limits its own edge cut-off to $W/2$
``as this fills the window completely''. A transit whose mask joins a data gap into a break lies between two segments
and is not tested.

\subsection{Counting rule} \label{app:counting}

A predicted transit is counted only if all four conditions hold:
\begin{enumerate}[label=(\roman*)]
\item at least 80\,\% of its $\Tfourteen$ window has cadences after cleaning;
\item it has at least 2 cadences on each side between 0.75 and 3\,$\Tfourteen$ from the predicted mid-time;
\item its whole $\Tfourteen$ window lies inside one segment and outside the edge zones of that segment;
\item in Sector 110, it is one of the frozen epochs of Tables~\ref{tab:timesA} and \ref{tab:timesB}.
\end{enumerate}
A transit that fails any of these conditions is not tested. It is listed with its status and depth, never enters the
statistic, and is never counted as a missing transit.


\subsection{Notes on the expected depths} \label{app:depthnotes}

These notes explain the values of Table~\ref{tab:depths}.
\begin{enumerate}
\item In Sector 32 the two B events fall about 2.0 and 2.3\,h after transits of A. The masks of the two signals then
  join into a hole longer than the 0.3\,d break tolerance, and the detrender fits the two sides separately. By the
  edge-zone rule (Appendix~\ref{app:edges}) these transits are not tested. In Sector 110 the closest pairs (A116 and
  B33, A109 and B31) are 7.4 and 7.7\,h apart, and their masks are separated by 2.1 and 2.4\,h of data. The same
  happens there only if most of that data is lost (about 1.6\,h next to the A mask, about 2.2\,h next to the B mask),
  and then both events of the pair are not tested.
\item The Sector 6 value is measured at 30-minute cadence, where the box statistic is lower than at 2-minute cadence. A
  transit-model check made for draft v2 (kept by me and not part of this deposit, so not verifiable from it; given for information only) gives a ratio of box depth to
  maximum depth of about 0.80 at 30\,min against 0.85 to 0.88 at 2\,min. One of the four counted Sector 6 transits of A
  (epoch $-793$) passes the 80\,\% coverage rule with 0.802, and in Sector 98 epoch 10 also passes with 0.802. The lower box
  depth at 30-minute cadence and these two marginal coverages are accepted as they are.
\end{enumerate}
